\documentclass[../../../include/open-logic-section]{subfiles}

\begin{document}

\olfileid{sth}{ordinals}{ordtype}
\olsection{ક્રમપ્રકારો તરીકે ક્રમવાચક સંખ્યાઓ}

હવે આપણી પાસે પ્રતિસ્થાપન છે અને આપણે $\ZFminus$માં કામ કરીએ છીએ.
તેથી અંતે આપણે ઇચ્છેલું પરિણામ સાબિત કરી શકીએ:

\begin{thm}\ollabel{thmOrdinalRepresentation}
દરેક સુક્રમ એક અને માત્ર એક ક્રમવાચક સંખ્યા સાથે ક્રમ-એકરૂપ છે.
\end{thm}

\begin{proof}
ધારો કે $\tuple{A, <}$ સુક્રમ છે.
\olref[basic]{ordisoidentity} મુજબ તે વધુમાં વધુ એક ક્રમવાચક
સંખ્યા સાથે ક્રમ-એકરૂપ છે. તેથી વિરોધાભાસ માટે ધારો કે
$\tuple{A, <}$ \emph{કોઈ પણ} ક્રમવાચક સંખ્યા સાથે
ક્રમ-એકરૂપ નથી. પહેલાં આપણે ``$\tuple{A, <}$ને બને તેટલું નાનું
બનાવીશું''. વિગતવાર: કોઈ ઉચિત આરંભિક ખંડ $\tuple{A_a,
<_a}$ કોઈ ક્રમવાચક સંખ્યા સાથે ક્રમ-એકરૂપ ન હોય, તો આ ગુણધર્મ
ધરાવતું લઘુતમ $a \in A$ મળે; પછી $B = A_a$ અને $\mathord{\lessdot} =
\mathord{<_a}$ લો. અન્યથા $B = A$ અને $\mathord{\lessdot} =
\mathord{<}$ લો.

રચના મુજબ $B$નો દરેક ઉચિત આરંભિક ખંડ કોઈ ક્રમવાચક સંખ્યા સાથે
ક્રમ-એકરૂપ છે, અને ઉપર મુજબ એ સંખ્યા અનન્ય છે. તેથી
પ્રતિસ્થાપનથી નીચેનો ગણ અસ્તિત્વમાં છે અને તે વિધેય છે:
\[
	f = \Setabs{\tuple{\beta, b}}{b \in B\text{ અને }\ordeq{\beta}{\tuple{B_b, \lessdot_b}}}
\]
વિરોધાભાસ પૂર્ણ કરવા આપણે બતાવીશું કે કોઈ ક્રમવાચક સંખ્યા
$\alpha$ માટે $f$ એ $\alpha \to B$ એકરૂપતા છે.

સ્પષ્ટ છે કે $\ran{f}
= B$. વળી \olref[iso]{lemordsegments} મુજબ $f$ ક્રમ જાળવે છે,
એટલે કે $\gamma \in \beta$ જો અને માત્ર જો
$f(\gamma) \lessdot f(\beta)$. હવે $\dom{f}$ ક્રમવાચક
સંખ્યા છે તે બતાવવા \olref[basic]{corordtransitiveord}
મુજબ $\dom{f}$ પરંપરિત છે એટલું બતાવવું પૂરતું છે. તેથી
$\beta \in \dom{f}$ લો, એટલે કે કોઈ $b$ માટે
$\ordeq{\beta}{\tuple{B_b, \lessdot_b}}$. જો $\gamma \in
\beta$ હોય, તો \olref[iso]{wellordinitialsegment} મુજબ
$\gamma \in \dom{f}$. સામાન્યકરણ કરતાં $\beta \subseteq
\dom{f}$.
\end{proof}

આ પરિણામથી આપણે \olref[vn]{sec}થી ઇચ્છતા હતા તે વ્યાખ્યા આપી
શકીએ છીએ:

\begin{defn}
જો $\tuple{A, <}$ સુક્રમ હોય, તો તેનો ક્રમપ્રકાર
$\ordtype{A, <}$ એવી અનન્ય ક્રમવાચક સંખ્યા $\alpha$ છે કે
$\ordeq{\tuple{A, < }}{\alpha}$.
\end{defn}

વળી, આ વ્યાખ્યાથી બે સરસ સિદ્ધાંતો મળે છે:

\begin{cor}\ollabel{ordtypesworklikeyouwant}
$\tuple{A, <}$ અને $\tuple{B, \lessdot}$ સુક્રમો હોય ત્યારે:
\begin{align*}
	\ordtype{A, <} = \ordtype{B, \lessdot}&\text{ જો અને માત્ર જો }\ordeq{\tuple{A, <}}{\tuple{B, \lessdot}}\\
	\ordtype{A, <} \in \ordtype{B, \lessdot}&\text{ જો અને માત્ર જો }\ordeq{\tuple{A, <}}{\tuple{B_b, \lessdot_b}}\text{ કોઈ }b \in B\text{ માટે}
\end{align*}
\end{cor}

\begin{proof}
સમાનતાનું નિવેદન \olref[basic]{ordisoidentity}થી મળે છે. બીજું
નિવેદન સાબિત કરવા $\ordtype{A, <} = \alpha$ અને $\ordtype{B,
\lessdot} = \beta$ લો, અને $f \colon \beta \to \tuple {B, \lessdot}$
આપણી એકરૂપતા હોય. તો:
\begin{align*}
	\alpha \in \beta&\text{ જો અને માત્ર જો }\funrestrictionto{f}{\alpha} \colon \alpha \to B_{f(\alpha)}\text{ એકરૂપતા છે}\\
	&\text{ જો અને માત્ર જો }\ordeq{\tuple{A, <}}{\tuple{B_{f(\alpha)}, \lessdot_{f(\alpha)}}}\\
	&\text{ જો અને માત્ર જો }\ordeq{\tuple{A, <}}{\tuple{B_b, \lessdot_b}}\text{ કોઈ $b \in B$ માટે}
\end{align*}
\sourcecorrection{OLMD-117}{મૂળના પ્રદર્શનની પહેલી હરોળમાં
$f(\alpha)$ ફક્ત $\alpha \in \beta$ હોય ત્યારે જ
વ્યાખ્યાયિત છે. તેથી તેની જમણી બાજુને આ શરત સાથે વાંચવી;
નહિતર એ સમકક્ષતાનું જમણું વિધાન વ્યાખ્યાયિત જ નથી.
તે જ રીતે $f \colon \beta \to \tuple {B, \lessdot}$માં
વિધેયનું લક્ષ્ય-ગણ $B$ છે અને $\tuple{B, \lessdot}$ એ
લક્ષ્યની ક્રમરચના દર્શાવે છે.}
\olref[iso]{ordisounique},
\olref[iso]{wellordinitialsegment}, અને
\olref[basic]{ordissetofsmallerord} મુજબ.
\end{proof}

\end{document}
